%=============================================================================!
% 15.0  CHAPTER OPENING                                                       !
%=============================================================================!
Many results of a simulation are averages over a run of finite length.
A second run from another start would give a slightly different value.
Chapters~11 to~14 quoted such averages, sometimes with the spread of a
few runs, but left open how accurately they were known. An error bar
makes that question explicit when we compare a result with experiment,
another code or another thermostat. It does not remove a bias: an early
interval in which the system is still settling shifts the average, and
extending the run only dilutes that shift.

We first find the error of a mean from independent samples, then account
for the memory between successive samples of a trajectory. Block averages
provide a second estimate and a visible check on whether the record is
long enough. These tools let us choose where to begin averaging, test for
drift and estimate the time needed for a chosen precision.

Convergence in time is only part of the problem. We also test the size of
the periodic box and whether the thermostat samples its stated ensemble.
The chapter ends with a protocol for preparing a run and a numerical
criterion for each reported quantity. Chapter~16 uses these checks to
measure structure and motion.
